\documentclass[aspectratio=169]{beamer}

\usetheme{Madrid}
\usepackage{booktabs}
\usepackage{xcolor}
\setbeamertemplate{navigation symbols}{}

\title[An ownable, shareable LULC classifier]{Extending the LULC classifier: user control, merging, and sharing}
\author[Salil Gujar]{Salil Gujar \\ \footnotesize M.Tech CSE \quad Entry No.\ 2025MCS2106}
\date{June 2026}

\begin{document}

\frame{\titlepage}

% ------------------------------------------------------------------
\begin{frame}{Where we are, and what was advised}
\begin{columns}[T]
\begin{column}{0.5\textwidth}
\emph{Done so far}
\begin{itemize}\small
  \item A 4-class base map over India at 10\,m, plus a living hierarchy: split or add a class,
  give a few example polygons, and retrain that node on the fly.
  \item A running, git-backed zoo of Model and Dataset cards: browse them, apply a model to the
  map, and publish to a shared repository, with metrics on each card.
\end{itemize}
\end{column}
\begin{column}{0.5\textwidth}
\emph{What was advised}
\begin{itemize}\small
  \item Let users own and share their scheme: choose the base classes, choose the inference
  data, and save and reload what they built.
  \item Add a merge operation: relabelling classes across models.
  \item Recommend where a model applies.
  \item On publish, store the model, record the contributor, and link the data.
\end{itemize}
\end{column}
\end{columns}
\end{frame}

% ------------------------------------------------------------------
\begin{frame}{Choosing the starting base classes}
\begin{itemize}
  \item On arrival, the user chooses a starting scheme rather than being fixed to one:
  \begin{itemize}\small
    \item IndiaSAT: greenery, water, built-up, barren, the calibrated default.
    \item WorldCover base (7 classes): tree, shrubland, grassland, cropland, built-up, bare, water.
  \end{itemize}
  \item The WorldCover base is trained from WorldCover-labelled Alpha Earth points, keeping
  only the classes with real coverage over India; the very rare ones are dropped for lack of
  support.
  \item It is an alternate starting point, not a more accurate one: WorldCover labels are weak,
  so it trails IndiaSAT and is presented that way.
  \item Both bases appear as cards in the zoo, so selecting a base scheme is a first-class
  action.
\end{itemize}
\end{frame}

% ------------------------------------------------------------------
\begin{frame}{Choosing the inference data}
\begin{itemize}
  \item The same trained model can classify a different year. Alpha Earth's embeddings are
  temporally consistent, so the model is reused and features are simply sampled at the chosen
  year.
  \item Realistic mode offers any year from 2017 to 2024. Detailed mode stays at 2024, the only
  year Tessera covers over India.
  \item No retraining is involved, and the chosen year is recorded so a result is reproducible.
  \item In practice, the same area in 2022 and 2024 shows a genuine difference in the class mix.
\end{itemize}
\end{frame}

% ------------------------------------------------------------------
\begin{frame}{Merging classes across models}
\begin{itemize}
  \item A merge is, in essence, a relabelling: take a class produced at one node and a class
  produced at another, possibly from different models, and declare them one new class.
  \item It runs as a correction layer applied after the classifiers, not as retraining.
  \item For example, relabel tea, from the greenery split, together with mining, from the
  barren split, into a single extractive class; the map and the counts reflect it at once.
  \item Merges are stored as rules, so the result is reproducible and a merge can be removed.
\end{itemize}
\end{frame}

% ------------------------------------------------------------------
\begin{frame}{Saving and reloading a scheme}
\begin{itemize}
  \item Users combine models in different ways to reach their own classifier. They can now
  download the hierarchy they built, along with the ordered steps that built it, as a JSON
  file, and load it back later.
  \item There are no accounts or sessions: the JSON file is the save. This keeps the tool light
  while letting people resume exactly where they left off.
  \item On reload, the tree is validated, each split is rebound to its trained model, and any
  split that needs retraining is reported.
\end{itemize}
\end{frame}

% ------------------------------------------------------------------
\begin{frame}{Recording the order of operations}
\begin{itemize}
  \item The structure of a hierarchy does not, by itself, capture how it was reached. The order
  of operations matters: a later split or merge depends on what came before.
  \item Every tree-changing action, split, add, retrain, apply, merge, and base switch, is
  appended to an ordered operation log.
  \item The log travels inside the saved scheme, so a shared or restored hierarchy reproduces
  the same output, and the workflow that produced a unique classifier is preserved.
\end{itemize}
\end{frame}

% ------------------------------------------------------------------
\begin{frame}{Recommending where a model applies}
\begin{itemize}
  \item Each model card now suggests where it fits: apply after a given class, derived from the
  model's node in the hierarchy and any WorldCover mapping its classes carry.
  \item The suggestion is also aware of the area in view, indicating whether the model is valid
  for the region currently on screen.
  \item It is computed from existing card metadata rather than hand-curated, so it stays correct
  as the zoo grows.
\end{itemize}
\end{frame}

% ------------------------------------------------------------------
\begin{frame}{Measuring data spread at a chosen scale}
\begin{itemize}
  \item Training data that is clustered in one area tends to skew a model. We summarise this as
  a spatial-diversity measure: how evenly the labelled polygons are spread.
  \item The grid cell on which the spread is measured is now chosen by the user, rather than
  fixed, so the measure can be read at the scale that matters for a given dataset.
  \item A finer grid reveals more occupied cells and usually a higher spread; the value updates
  immediately on the dataset card.
\end{itemize}
\end{frame}

% ------------------------------------------------------------------
\begin{frame}{Publishing: storing the model}
\begin{itemize}
  \item When a user publishes, the trained model file is committed into the shared zoo
  repository. The models are small, so a published model is usable by anyone who pulls it,
  not merely described.
  \item The repository continues to hold only cards and these model files; large training
  tables and tiles are kept out.
  \item Publishing also records the contributor, a GitHub handle or an email, on the card, so
  it is clear who shared each model.
\end{itemize}
\end{frame}

% ------------------------------------------------------------------
\begin{frame}{Publishing: the data and the contract}
\begin{itemize}
  \item For each training dataset, publishing asks, at that moment, for a public link to where
  the data came from, and shows it on the card.
  \item A private upload with no public link is never shared: it stays local and is used only to
  train the model. The labelled data remains the user's own.
  \item This keeps the published catalogue rich in provenance while honouring the agreement that
  the data itself belongs to whoever uploaded it.
\end{itemize}
\end{frame}

% ------------------------------------------------------------------
\begin{frame}
\centering
\Large Thank you
\end{frame}

\end{document}
